% !TEX program = pdflatex
\documentclass[11pt,aspectratio=169]{beamer}

% Основной движок — pdfLaTeX: Latin Modern с кодировкой T2A есть
% в любом дистрибутиве, поэтому кириллица собирается и в Overleaf,
% и локально без установки шрифтов.
%
% Ветка XeLaTeX оставлена на случай, когда системные шрифты доступны
% и хочется вид посимпатичнее. Осторожно: при отсутствии указанного
% шрифта XeLaTeX не падает, а молча рисует пустые страницы.
\usepackage{iftex}
\ifPDFTeX
  \usepackage[T2A]{fontenc}
  \usepackage[utf8]{inputenc}
  \usepackage[english,russian]{babel}
  \usepackage{lmodern}
\else
  \usepackage{fontspec}
  \setsansfont{DejaVu Sans}[Scale=0.92]
  \setmonofont{DejaVu Sans Mono}[Scale=0.80]
\fi
\usepackage{amsmath,amssymb}
\usepackage{booktabs}
\usepackage{xcolor}

\definecolor{acc}{HTML}{7A2E1E}
\definecolor{dim}{HTML}{555555}
\definecolor{boxbg}{HTML}{F4F1EC}

\usetheme{default}
\usecolortheme{default}
\setbeamertemplate{navigation symbols}{}
\setbeamercolor{frametitle}{fg=acc}
\setbeamerfont{frametitle}{size=\large,series=\bfseries}
\setbeamercolor{structure}{fg=acc}
\setbeamertemplate{itemize item}{\textbullet}
\setbeamertemplate{itemize subitem}{--}
\setbeamerfont{footline}{size=\tiny}
\setbeamertemplate{footline}{%
  \hfill\usebeamerfont{footline}\color{dim}\insertframenumber\hspace{2em}\vspace{1em}}

\newcommand{\key}[1]{%
  \begin{center}\colorbox{boxbg}{\parbox{0.92\textwidth}{\vspace{2pt}#1\vspace{2pt}}}\end{center}}
\newcommand{\aside}[1]{{\color{dim}\small #1}}
\newcommand{\R}{\mathbb{R}}
\newcommand{\E}{\mathbb{E}}
\newcommand{\Var}{\operatorname{Var}}
\newcommand{\KL}{\operatorname{KL}}

\title{\bfseries Дообучение и выравнивание}
\subtitle{SFT, LoRA и предпочтения: математика и практика}
\date{}

\begin{document}

\begin{frame}[plain]
\titlepage
\vspace{-2em}
\aside{vlm-recipes}
\end{frame}

% ══════════════════════════════════════════════════════════════
\section{Постановка}

\begin{frame}{Две разные задачи}

Дообучение под свою предметную область распадается надвое.

\vspace{1em}
\begin{columns}[t]
\column{0.47\textwidth}
\textbf{\color{acc}Научить форме}

Воспроизводить эталонный ответ: стиль, структура, дисциплина вызова
инструментов.

\vspace{0.5em}
\aside{Это SFT.}

\column{0.47\textwidth}
\textbf{\color{acc}Научить различать}

Предпочитать один хороший ответ другому хорошему.

\vspace{0.5em}
\aside{Это выравнивание.}
\end{columns}

\vspace{1.5em}
\key{Второе не следует из первого. Разберём почему — и что с этим делать.}

\end{frame}

\begin{frame}{Что модель делает на самом деле}

Функция одного вида: последовательность токенов на входе, распределение
над следующим токеном на выходе.
\[
p_\theta(\,\cdot \mid x_1,\dots,x_{t-1}) \in \Delta^{|V|}
\]

\vspace{1em}
Больше ничего. Диалог, вызов инструментов, следование инструкции —
надстройка над этим единственным действием.

\vspace{1em}
\key{Всё дообучение сводится к двум вопросам: \textbf{на каких токенах
считается градиент} и \textbf{в каком подпространстве разрешено двигать
веса}.}

\end{frame}

% ══════════════════════════════════════════════════════════════
\section{SFT}

\begin{frame}{Что оптимизируется}

Отрицательный логарифм правдоподобия — кросс-энтропия по всей
последовательности:
\[
\mathcal{L}(\theta) = -\sum_{t=1}^{T} \log p_\theta(x_t \mid x_{<t})
\]

\vspace{0.5em}
Проблема: сумма идёт по \emph{всем} позициям — включая вопрос
пользователя и результаты вызовов инструментов.

\vspace{1em}
\key{Модель, обученная так, учится \textbf{порождать вопросы
пользователя} и \textbf{выдумывать вывод инструментов}.}

\end{frame}

\begin{frame}{Маскирование}

Сумма ограничивается множеством $A$ позиций реплик ассистента:
\[
\mathcal{L}_{\text{SFT}}(\theta) = -\frac{1}{|A|}\sum_{t \in A}
\log p_\theta(x_t \mid x_{<t})
\]

\vspace{1em}
Существенно: $x_{<t}$ по-прежнему включает \emph{весь} контекст.
Маскирование убирает позиции из потерь, но не из обусловливания.

\vspace{0.5em}
Технически — значение $-100$ в метках, которое \texttt{cross\_entropy}
пропускает.

\vspace{1em}
\aside{Диагностика: если доля обучаемых токенов $|A|/T$ измеряется
единицами процентов, почти весь бюджет батча уходит на проход без
градиента.}

\end{frame}

\begin{frame}{Почему нельзя дообучать всё}

Для AdamW со смешанной точностью, $P$ параметров:

\vspace{0.5em}
\begin{center}
\begin{tabular}{@{}lr@{}}
\toprule
веса bf16 & $2P$ \\
мастер-копия fp32 & $4P$ \\
градиенты & $2P$ \\
моменты $m,v$ & $8P$ \\
\midrule
\textbf{итого} & $\approx 16P$ байт \\
\bottomrule
\end{tabular}
\end{center}

\vspace{1em}
Для $P = 9\cdot 10^9$ — около \textbf{144 ГБ}, без учёта активаций.

\vspace{0.5em}
\key{Все методы ниже атакуют разные слагаемые этой суммы.}

\end{frame}

\begin{frame}{LoRA: конструкция}

Наблюдение: приращение $\Delta W$, нужное для адаптации к узкой задаче,
хорошо приближается матрицей малого ранга.

\[
\Delta W = BA, \qquad B \in \R^{d\times r},\; A \in \R^{r\times k},
\qquad r \ll \min(d,k)
\]
\[
h = W_0 x + \gamma\, BAx, \qquad \gamma = \frac{\alpha}{r}
\]

\vspace{0.5em}
Параметров: $r(d+k)$ вместо $dk$. При $d=k=4096$, $r=32$ — в 64 раза меньше.

\vspace{1em}
\key{Замороженные $W_0$ не требуют ни градиентов, ни моментов. Три
слагаемых из четырёх исчезают.}

\end{frame}

\begin{frame}{Инициализация: почему $B = 0$}

$A$ случайное, $B$ нулевое $\;\Rightarrow\; \Delta W = 0$ на старте:
модель выходит ровно из претрейна.

\vspace{0.8em}
Градиенты, при $g = \partial L/\partial h$:
\[
\frac{\partial L}{\partial B} = \gamma\, g\,(Ax)^\top,
\qquad
\frac{\partial L}{\partial A} = \gamma\, B^\top g\, x^\top
\]

\vspace{0.5em}
Обнуляя одну матрицу, вы убиваете градиент по \emph{другой}. Обе
односторонние схемы обучаются — зеркально друг другу.

\vspace{0.8em}
\key{Сломана только двусторонняя: обе нулевые $\Rightarrow$ оба
градиента нулевые навсегда.}

\end{frame}

\begin{frame}{Масштаб: почему $\alpha/r$ плох при больших рангах}

Элемент произведения — сумма $r$ слагаемых:
\[
(BA)_{ij} = \sum_{s=1}^{r} B_{is}A_{sj},
\qquad
\Var\big[(BA)_{ij}\big] \approx r\,\sigma_B^2\sigma_A^2
\]

Среднеквадратичная величина растёт как $\sqrt{r}$ — случайное блуждание,
знаки частично гасятся. С множителем $\gamma = \alpha/r$:
\[
\gamma\cdot\sqrt{r}\,\sigma_B\sigma_A
= \frac{\alpha\,\sigma_B\sigma_A}{\sqrt{r}}
\;\xrightarrow[r\to\infty]{}\; 0
\]

\key{Вклад адаптера \textbf{затухает}. Поднимая ранг, вы одновременно
глушите адаптер. Правильный масштаб — $\gamma = \alpha/\sqrt{r}$,
это \textbf{rsLoRA}.}

\end{frame}

\begin{frame}{DoRA: разделение величины и направления}

При полном дообучении норма столбца и его направление меняются
\emph{независимо}. LoRA такой свободы не даёт: одно приращение $BA$
управляет обоими.

\vspace{0.8em}
\[
W = m \odot \frac{V}{\lVert V\rVert_c}
\qquad\Longrightarrow\qquad
W' = m \odot \frac{W_0 + BA}{\lVert W_0 + BA\rVert_c}
\]

\vspace{0.5em}
Обучаются $m$, $A$, $B$. Нормировка развязывает две степени свободы.

\vspace{1em}
\aside{На низких рангах ($r \le 16$) устойчиво обходит LoRA. Цена —
вычисление нормы на каждом шаге, около 20\% скорости.}

\end{frame}

\begin{frame}{QLoRA: квантование базы}

Веса сети распределены приблизительно нормально — равномерная сетка
int4 тратит разрешение на хвосты.

\vspace{0.5em}
\textbf{NF4} строит уровни как квантили нормального распределения:
\[
q_i \;\propto\; Q_{\mathcal{N}}\!\left(\frac{i+0.5}{2^b}\right)
\]

Блоки по 64 со своим масштабом; двойное квантование сжимает и масштабы
до $\approx 0{,}127$ бита на параметр.

\vspace{0.8em}
\[
h = \text{dequant}_{\text{bf16}}(W^{\text{NF4}})\,x + \gamma BAx
\]

\key{Умножение всё равно в bf16 — \textbf{QLoRA экономит память, но не
ускоряет}. Брать, только когда bf16-база не помещается.}

\end{frame}

\begin{frame}{Сводка: что выбрать}

\begin{center}
\small
\begin{tabular}{@{}lp{4.2cm}p{4.6cm}@{}}
\toprule
\textbf{Метод} & \textbf{Когда} & \textbf{Цена} \\
\midrule
full & максимум качества, ресурсы есть & $16P$ байт памяти \\
lora & базовый выбор & $r$ выше 32 бесполезен \\
\textbf{rslora} & \textbf{всегда вместо lora} & нет \\
dora & низкие ранги, тонкий стиль & $-20\%$ скорости \\
qlora & база не влезает в bf16 & медленнее lora \\
galore & полноранговое при ограниченной памяти & SVD раз в $N$ шагов \\
\bottomrule
\end{tabular}
\end{center}

\vspace{1em}
\aside{Передний край: LoRA+ (разные lr для $A$ и $B$), PiSSA
(инициализация главными компонентами $W_0$), VeRA (общие замороженные
$A,B$), Spectrum и LISA (выбор подмножества слоёв).}

\end{frame}

% ══════════════════════════════════════════════════════════════
\section{Мотивация выравнивания}

\begin{frame}{Чего не умеет SFT}

SFT максимизирует правдоподобие \emph{эталонного} ответа. Это учит
воспроизводить хорошие ответы.

\vspace{1em}
Но правильных ответов на один вопрос много, и различаются они оттенком.

\vspace{1em}
\key{SFT видит \textbf{только положительные примеры} и никогда не
наблюдает, чего говорить \emph{не} следует. Градиент поднимает
вероятность эталона, но \textbf{не опускает} вероятность правдоподобной
альтернативы.}

\end{frame}

\begin{frame}{Пример: оба ответа грамотные}

\small
Запрос: \textit{«Сформулируй мне гипотезу для диплома»}

\vspace{0.8em}
\begin{columns}[t]
\column{0.47\textwidth}
\textbf{\color{acc}Выбранный}

\vspace{0.3em}
\colorbox{boxbg}{\parbox{0.95\textwidth}{\vspace{2pt}\footnotesize
Гипотеза вырастает из проблемы, а её в работе пока нет. Какое
противоречие вы заметили?\vspace{2pt}}}

\column{0.47\textwidth}
\textbf{\color{acc}Отвергнутый}

\vspace{0.3em}
\colorbox{boxbg}{\parbox{0.95\textwidth}{\vspace{2pt}\footnotesize
Гипотеза: внедрение дистанционного формата приводит к снижению
успеваемости студентов первого курса.\vspace{2pt}}}
\end{columns}

\vspace{1.2em}
\normalsize
Второй ответ грамотен, уместен и звучит компетентно. Он нарушает одно:
содержательное решение принимает студент.

\vspace{0.8em}
\key{Для SFT эти два ответа неразличимы: он видит только первый и не
знает, что второй существует.}

\end{frame}

% ══════════════════════════════════════════════════════════════
\section{Выравнивание}

\begin{frame}{Классический путь: RLHF}

По парам $(y_w \succ y_l)$ обучается модель награды. Вероятность
предпочтения — модель Брэдли–Терри:
\[
p(y_w \succ y_l \mid x) = \sigma\big(r(x,y_w) - r(x,y_l)\big)
\]

Затем политика оптимизируется при KL-штрафе:
\[
\max_{\pi}\;
\E_{x\sim D,\,y\sim\pi}\big[r(x,y)\big]
\;-\;\beta\,\KL\!\big(\pi(\cdot\mid x)\,\|\,\pi_{\text{ref}}(\cdot\mid x)\big)
\]

\vspace{0.5em}
Штраф необходим: без него политика уходит в вырожденные тексты.

\vspace{0.5em}
\aside{Цена: четыре модели в памяти (политика, опорная, награды, критик)
и капризная настройка.}

\end{frame}

\begin{frame}{DPO, шаг 1: оптимальная политика}

Задача выше решается \emph{аналитически}. Для фиксированного $x$
решение больцмановского вида:

\vspace{0.5em}
\[
\pi^*(y\mid x) = \frac{1}{Z(x)}\,
\pi_{\text{ref}}(y\mid x)\,\exp\!\Big(\tfrac{1}{\beta}\,r(x,y)\Big)
\]
\[
Z(x) = \sum_{y}\pi_{\text{ref}}(y\mid x)\,e^{r(x,y)/\beta}
\]

\vspace{1em}
Нормировочная сумма $Z(x)$ берётся по всем возможным текстам — посчитать
её нельзя. Пока это выглядит тупиком.

\end{frame}

\begin{frame}{DPO, шаг 2: обращение}

Логарифмируем и выражаем награду через политику:

\vspace{0.5em}
\[
r(x,y) = \beta\,\log\frac{\pi^*(y\mid x)}{\pi_{\text{ref}}(y\mid x)}
+ \beta\,\log Z(x)
\]

\vspace{1.5em}
\key{Награда оказалась \textbf{выражена через саму политику}. Отдельная
модель награды не нужна: любая политика неявно задаёт награду,
и наоборот.}

\end{frame}

\begin{frame}{DPO, шаг 3: подстановка}

Подставляем в Брэдли–Терри. Слагаемое $\beta\log Z(x)$ одинаково для
$y_w$ и $y_l$ — \textbf{сокращается в разности} вместе с неберущейся суммой:

\vspace{0.5em}
\[
p(y_w \succ y_l \mid x) = \sigma\!\left(
\beta\log\frac{\pi_\theta(y_w\mid x)}{\pi_{\text{ref}}(y_w\mid x)}
-\beta\log\frac{\pi_\theta(y_l\mid x)}{\pi_{\text{ref}}(y_l\mid x)}
\right)
\]

\vspace{0.5em}
Максимизация правдоподобия предпочтений:
\[
\mathcal{L}_{\text{DPO}} = -\,\E
\left[\log\sigma\!\left(
\beta\log\frac{\pi_\theta(y_w\mid x)}{\pi_{\text{ref}}(y_w\mid x)}
-\beta\log\frac{\pi_\theta(y_l\mid x)}{\pi_{\text{ref}}(y_l\mid x)}
\right)\right]
\]

\end{frame}

\begin{frame}{Что делает градиент DPO}

Обозначим неявный отступ
$\hat\Delta = \hat r_\theta(x,y_w) - \hat r_\theta(x,y_l)$:

\vspace{0.5em}
\[
\nabla_\theta \mathcal{L}_{\text{DPO}} = -\beta\,
\underbrace{\sigma(-\hat\Delta)}_{\text{вес примера}}
\Big[\nabla_\theta\log\pi_\theta(y_w\mid x)
- \nabla_\theta\log\pi_\theta(y_l\mid x)\Big]
\]

\vspace{1em}
Вес $\sigma(-\hat\Delta)$ велик там, где модель \emph{ошибается}
в предпочтении, и близок к нулю там, где она уже права.

\vspace{0.8em}
\key{Обучение само концентрируется на трудных парах. Отсюда
устойчивость метода и малое требуемое число примеров.}

\end{frame}

\begin{frame}{ORPO: один этап вместо двух}

DPO требует предварительного SFT и опорной модели в памяти. ORPO
обходится без обоих.

\vspace{0.8em}
Шансы последовательности и их отношение:
\[
\text{odds}_\theta(y\mid x) = \frac{\pi_\theta(y\mid x)}{1-\pi_\theta(y\mid x)},
\qquad
\text{OR} = \frac{\text{odds}_\theta(y_w\mid x)}{\text{odds}_\theta(y_l\mid x)}
\]
\[
\mathcal{L}_{\text{ORPO}} = \mathcal{L}_{\text{SFT}}
- \lambda\,\log\sigma\big(\log \text{OR}\big)
\]

\vspace{0.5em}
Первое слагаемое тянет к эталону, второе разводит выбранный
и отвергнутый.

\vspace{0.5em}
\aside{Отношение шансов выбрано вместо отношения вероятностей намеренно:
наказывает мягче и не разрушает усвоенное на SFT.}

\end{frame}

\begin{frame}{SimPO: без опорной модели}

Неявная награда DPO — \emph{сумма} по токенам, а генерация
максимизирует \emph{среднюю} логвероятность. Рассогласование:
DPO систематически предпочитает длинные ответы.

\vspace{0.8em}
\[
r_{\text{SimPO}}(x,y) = \frac{\beta}{|y|}\log\pi_\theta(y\mid x),
\qquad
\mathcal{L} = -\log\sigma\big(r(x,y_w) - r(x,y_l) - \gamma\big)
\]

\vspace{0.5em}
Награда нормирована на длину, опорная модель не нужна.

\vspace{0.5em}
Порог $\gamma > 0$ требует превосходства \emph{с запасом}.

\vspace{0.5em}
\aside{Побочный эффект — метод не раздувает длину ответов.}

\end{frame}

\begin{frame}{KTO: конструкция}

Неявная награда — та же, что в DPO:
\[
r_\theta(x,y) = \log\frac{\pi_\theta(y\mid x)}{\pi_{\text{ref}}(y\mid x)}
\]

Вторую половину пары заменяет \textbf{точка отсчёта} — среднее отклонение
политики от опорной, оцениваемое по батчу на перемешанных парах:
\[
z_0 = \KL\big(\pi_\theta(y'\mid x)\,\|\,\pi_{\text{ref}}(y'\mid x)\big)
\]

\vspace{0.3em}
\aside{$z_0$ считается как константа и в градиенте не участвует.}

\end{frame}

\begin{frame}{KTO: функция ценности}

Ценность ответа измеряется \emph{относительно} $z_0$, и знак сравнения
зависит от метки:
\[
v(x,y) =
\begin{cases}
\lambda_D\,\sigma\big(\beta\,(r_\theta(x,y) - z_0)\big), & y \text{ помечен «хорошо»}\\[0.4em]
\lambda_U\,\sigma\big(\beta\,(z_0 - r_\theta(x,y))\big), & y \text{ помечен «плохо»}
\end{cases}
\]

\vspace{0.5em}
\[
\mathcal{L}_{\text{KTO}} = \E_{x,y\sim D}\big[\lambda_y - v(x,y)\big]
\]

\vspace{0.8em}
Для хорошего ответа награда должна \emph{превышать} $z_0$, для плохого —
\emph{отставать} от неё. Сигмоида играет роль функции ценности из теории
перспектив: вогнута на приобретениях, выпукла на потерях.

\end{frame}

\begin{frame}{KTO: почему пара не нужна}

\begin{columns}[t]
\column{0.46\textwidth}
\textbf{\color{acc}DPO}

Ответ сравнивается с \textbf{конкретным} отвергнутым:
\[
\hat r(y_w) - \hat r(y_l)
\]

\column{0.46\textwidth}
\textbf{\color{acc}KTO}

Ответ сравнивается со \textbf{средним} поведением модели:
\[
r_\theta(x,y) - z_0
\]
\end{columns}

\vspace{1em}
\key{Точка отсчёта $z_0$ заменяет вторую половину пары. Сравнивать есть
с чем, даже когда про каждый пример известно только «хорошо» или
«плохо».}

\vspace{0.8em}
Асимметрия $\lambda_U > \lambda_D$ — из теории перспектив: потери
переживаются острее приобретений. Практически это ещё и рычаг под
несбалансированные данные — жалоб обычно меньше, чем лайков, и вес
компенсирует перекос.

\vspace{0.5em}
\aside{Отсюда главное следствие: данные можно собирать из продакшена.
Лайк или жалоба уже является меткой.}

\end{frame}

\begin{frame}{Сводка методов выравнивания}

\begin{center}
\begin{tabular}{@{}llcp{4.6cm}@{}}
\toprule
\textbf{Метод} & \textbf{Данные} & \textbf{Опорная} & \textbf{Когда брать} \\
\midrule
DPO & пары & да & есть SFT-модель и память на две копии \\
ORPO & пары & нет & один этап, ограниченные ресурсы \\
SimPO & пары & нет & важна длина ответа \\
KTO & метки & да & пары собрать дорого \\
\bottomrule
\end{tabular}
\end{center}

\vspace{1.5em}
\aside{В trl 1.x отдельные \texttt{ORPOTrainer} и \texttt{CPOTrainer}
слиты в \texttt{DPOTrainer} с разными значениями \texttt{loss\_type}.
Математика та же, поменялся способ вызова.}

\end{frame}

% ══════════════════════════════════════════════════════════════
\section{Практика}

\begin{frame}{Векторы управления: без обучения вовсе}

Абстрактные свойства ответа кодируются в активациях приблизительно
линейно. Направление — разность средних по контрастным множествам:
\[
v_\ell = \frac{1}{|P|}\sum_{p\in P} h_\ell(p)
- \frac{1}{|N|}\sum_{n\in N} h_\ell(n),
\qquad
h_\ell \leftarrow h_\ell + s\,\frac{v_\ell}{\lVert v_\ell\rVert}
\]

\vspace{0.5em}
Десятки пар и один проход — минуты вместо дней.

\vspace{0.8em}
\key{Если вектор не даёт эффекта ни при каком $s$, свойство не
выделяется линейно, и дообучение на тех же примерах скорее всего тоже
не поможет. \textbf{Самая дешёвая проверка достижимости.}}

\end{frame}

\begin{frame}{Лестница приёмов по стоимости}

\begin{center}
\small
\begin{tabular}{@{}llp{6.2cm}@{}}
\toprule
\textbf{Приём} & \textbf{Стоимость} & \textbf{Ограничение} \\
\midrule
системный промпт & 0 & перебивается текстом пользователя \\
префилл ответа & 0 & задаёт форму, не содержание \\
грамматика при декодировании & 0 & гарантирует формат, не факты \\
модель-судья & +1 вызов & удвоенная задержка \\
вектор управления & часы & грубый, сдвигает поведение целиком \\
SFT & дни & нужны эталонные ответы \\
DPO / ORPO & дни & нужны пары предпочтений \\
\bottomrule
\end{tabular}
\end{center}

\vspace{1em}
\key{Двигаться сверху вниз и останавливаться, как только результат
достигнут.}

\end{frame}

\begin{frame}{Что мерить}

Метрика всегда \textbf{пара}, и порознь обе бессмысленны:

\vspace{1em}
\begin{center}
\begin{tabular}{@{}lll@{}}
\toprule
\textbf{Метрика} & \textbf{На чём} & \textbf{Куда} \\
\midrule
доля срабатываний & целевые запросы & к 100\% \\
доля ложных срабатываний & контрольные запросы & к 0\% \\
\bottomrule
\end{tabular}
\end{center}

\vspace{1.2em}
Модель, срабатывающая всегда, даёт идеальную первую и катастрофическую
вторую.

\vspace{0.8em}
\key{Отсюда требование к данным: в каждом наборе нужна
\textbf{отрицательная группа} — примеры, похожие на целевые, но
требующие обычного ответа.}

\end{frame}

\begin{frame}{Итог}

\begin{enumerate}\setlength{\itemsep}{0.7em}
\item \textbf{SFT} учит форме. Ключевая деталь — маскирование: градиент
только по репликам ассистента.

\item \textbf{LoRA} делает это выполнимым. \texttt{rsLoRA} — дефолт,
масштаб $\alpha/\sqrt{r}$ вместо $\alpha/r$ снимает потолок по рангу.

\item \textbf{Выравнивание} учит различать. DPO выводится из RLHF
аналитически: награда выражается через саму политику, $Z(x)$
сокращается.

\item \textbf{Порядок работ:} приёмы без обучения $\to$ вектор
управления как проверка достижимости $\to$ SFT $\to$ выравнивание.
Каждый шаг может показать, что следующий не нужен.
\end{enumerate}

\end{frame}

\end{document}
